% !PRINTED 533
\weberpage{533}

Zunächst ist klar, dass eine Basis eines Functionals $\mu$ zugleich Basis
aller mit $\mu$ associirten Functionale ist, ferner dass jedes Element
$\alpha_1, \alpha_2, \ldots$ einer solchen Basis durch $\mu$ theilbar sein
muss, endlich dass aus einer Basis von $\mu$ alle anderen abgeleitet werden
können, wenn man eine lineare Substitution mit ganzen rationalen
Co\"efficienten und der Determinante $\pm 1$ anwendet.

Um nun eine Basis von $\mu$ zu bilden, gehen wir von einer Minimalbasis
(Basis von $\mathfrak{o}$) aus, die wir, wie oben, mit
\[
\omega_1, \omega_2, \ldots, \omega_n
\]
bezeichnen.

Da es positive ganze rationale Zahlen giebt, die durch $\mu$ theilbar sind,
so giebt es auch ganze rationale Zahlen $a$, für die das Product
$a\,\omega_1$ durch $\mu$ theilbar ist. Die \gesperrt{kleinste} positive unter diesen
Zahlen wollen wir mit $a_{1,1}$ bezeichnen und
\[
\alpha_1 = a_{1,1}\,\omega_1
\]
setzen. Sodann bezeichnen wir mit $a_{2,2}$ die \gesperrt{kleinste positive} ganze
rationale Zahl, für die sich ein ganzes rationales $a_{1,2}$ so bestimmen
lässt, dass
\[
\alpha_2 = a_{1,2}\,\omega_1 + a_{2,2}\,\omega_2
\]
durch $\mu$ theilbar wird, und fahren so fort. Wir bilden auf diese Weise
das System
\weq{3}{
\begin{aligned}
\alpha_1 &= a_{1,1}\,\omega_1 \\
\alpha_2 &= a_{1,2}\,\omega_1 + a_{2,2}\,\omega_2 \\
&\cdots\cdots\cdots\cdots\cdots\cdots\cdots\cdots\cdots\cdots \\
\alpha_n &= a_{1,n}\,\omega_1 + a_{2,n}\,\omega_2 + \cdots + a_{n,n}\,\omega_n,
\end{aligned}
}
\gesperrt{worin $a_{1,1}, a_{2,2}, \ldots, a_{n,n}$ die kleinsten positiven
ganzen rationalen Zahlen sind, die eine Bestimmung der ganzen rationalen
Zahlen $a_{1,2}, \ldots, a_{n-1,n}$ so ermöglichen, dass die ganzen Zahlen
$\alpha_1, \alpha_2, \ldots, \alpha_n$ durch $\mu$ theilbar werden.}

Jede Zahl $a_{r,r}$ ist hiernach unabhängig von den übrigen
dadurch definirt, dass sie die kleinste natürliche Zahl ist, für die sich
die zugehörigen ganzen rationalen Zahlen $a_{1,r}, a_{2,r}, \ldots,
a_{r-1,r}$ so bestimmen lassen, dass
\weq{4}{
\alpha_r = a_{1,r}\,\omega_1 + a_{2,r}\,\omega_2 + \cdots +
a_{r-1,r}\,\omega_{r-1} + a_{r,r}\,\omega_r
}
durch $\mu$ theilbar wird. In dem besonderen Falle, wo
$\omega_r$ selbst durch $\mu$ theilbar ist, hat man demnach $a_{r,r} = 1$
zu setzen, und
%CONT>>
